\subsection*{الف}
فرض کنید
$\omega_0=\frac{2\pi}{T}$.
این سیگنال را به صورت
$x(t)=-2+3p(t)$
می‌نویسیم که در آن
$p(t)$
یک پالس مربعی متناوب با دوره $T$ است و
\[
p(t)=
\begin{cases}
1 & 0\le t<t_1,\\
0 & t_1\le t<T.
\end{cases}
\]
پس ضرایب سری فوریه مختلط:
\[
c_n = -2\,\delta[n] + 3\,P_n,
\qquad
P_n=\frac{1}{T}\int_0^{t_1} e^{-jn\omega_0 t}\,dt.
\]
برای $n=0$:
\[
c_0=\frac{1}{T}\int_0^T x(t)\,dt=\frac{t_1-2(T-t_1)}{T}=\frac{3t_1}{T}-2.
\]
برای $n\neq 0$:
\begin{flalign*}
c_n
&=\frac{1}{T}\left(\int_0^{t_1} e^{-jn\omega_0 t}\,dt-2\int_{t_1}^{T} e^{-jn\omega_0 t}\,dt\right) \\
&=\frac{3}{j2\pi n}\left(1-e^{-jn\frac{2\pi}{T}t_1}\right)
=\frac{3}{\pi n}\,e^{-jn\pi\frac{t_1}{T}}\,\sin\!\left(n\pi\frac{t_1}{T}\right). &&
\end{flalign*}

\subsection*{ب}
اگر
$x_e(t)=\frac{x(t)+x(-t)}{2}$ و $x_o(t)=\frac{x(t)-x(-t)}{2}$،
آنگاه ضرایب سری فوریه مختلط آن‌ها:
\[
c_n^{(e)}=\frac{c_n+c_{-n}}{2},\qquad
c_n^{(o)}=\frac{c_n-c_{-n}}{2}.
\]
چون $x(t)$ حقیقی است، $c_{-n}=c_n^*$.
پس:
\[
c_0^{(e)}=c_0=\frac{3t_1}{T}-2,\qquad c_0^{(o)}=0,
\]
و برای $n\neq 0$:
\begin{flalign*}
c_n^{(e)}&=\Re\{c_n\}
=\frac{3}{2\pi n}\sin\!\left(2\pi n\frac{t_1}{T}\right),\\
c_n^{(o)}&=j\,\Im\{c_n\}
=-j\,\frac{3}{\pi n}\sin^2\!\left(n\pi\frac{t_1}{T}\right). &&
\end{flalign*}

\subsection*{ج}
طبق رابطه \lr{Parseval} برای سری فوریه مختلط:
\[
\sum_{n=-\infty}^{\infty}|c_n|^2=\frac{1}{T}\int_0^T |x(t)|^2\,dt
=\frac{1}{T}\left(\int_0^{t_1}1\,dt+\int_{t_1}^{T}4\,dt\right)
=4-\frac{3t_1}{T}.
\]
(همچنین $|c_n|^2=\frac{9}{\pi^2 n^2}\sin^2(n\pi\frac{t_1}{T})$ برای $n\neq 0$ و
$|c_0|^2=(\frac{3t_1}{T}-2)^2$.)